%%=============================================================================
%% Voorwoord
%%=============================================================================

\chapter*{\IfLanguageName{dutch}{Woord vooraf}{Preface}}%
\label{ch:voorwoord}

Deze bachelorproef vormt het sluitstuk van mijn opleiding Toegepaste Informatica aan HOGent. Doorheen mijn opleiding ontwikkelde ik een sterke interesse in bedrijfsprocessen en de manier waarop IT-systemen deze kunnen ondersteunen en optimaliseren. De keuze voor een onderwerp rond ERP-systemen en meer specifiek Odoo kwam dan ook voort uit de combinatie van technische implementatie en bedrijfskundige impact.

Tijdens deze bachelorproef kreeg ik de kans om een realistische case uit te werken rond Hanuman BV, een Belgische KMO actief in de verkoop en het onderhoud van koffieapparatuur. Het analyseren van de bestaande werkwijze en het uitwerken van een mogelijke ERP-oplossing bood mij de mogelijkheid om mijn theoretische kennis toe te passen in een praktijkgerichte context. Daarbij stond niet enkel de technische configuratie centraal, maar ook het begrijpen van de bedrijfsprocessen en het identificeren van verbeterpunten.

Het uitvoeren van deze bachelorproef was een leerrijk traject waarin ik zowel zelfstandig als gestructureerd heb leren werken. Daarnaast heb ik inzicht verworven in de uitdagingen die gepaard gaan met digitale transformatie binnen KMO’s, zoals datakwaliteit, change management en procesintegratie.

Graag wil ik mijn promotor, Antjon Tom, bedanken voor de begeleiding en feedback gedurende dit traject. Ook mijn co-promotor, Jurgen De Smet, wil ik bedanken voor zijn ondersteuning en inzichten. Tot slot wil ik mijn omgeving bedanken voor de motivatie en steun tijdens het uitwerken van deze bachelorproef.

\vspace{1cm}

Matisse De Smet